\begin{figure}[!t]
\centering
\includegraphics[width=\linewidth,height=0.72\textheight,keepaspectratio]{figures/Figure_1.pdf}
\caption{\textbf{Contribution dynamics without communication.} Panels show the six reasoning configurations at N=6 and R=0.40, 0.90, 0.98 and 1.10. Points show the ensemble mean group contribution in each round. Shading shows pointwise 95\% percentile intervals from 10,000 within-condition complete-trajectory bootstrap samples. All games use T=10, a fresh endowment E=20 each round, fixed membership, aggregate feedback and full history.}
\label{fig:1}
\end{figure}

\begin{figure}[!t]
\centering
\includegraphics[width=\linewidth,height=0.72\textheight,keepaspectratio]{figures/Figure_2.pdf}
\caption{\textbf{Incentive-dependent contribution at fixed group size.} All panels use N=6 and no communication. (a) Whole-run mean contribution over rounds 1--10 across all 11 scanned R values. (b) Final-round contribution at round 10 over the same scan. (c) Expanded view of the six social-dilemma R values from a. R=r/N is plotted at its actual numerical spacing. Vertical dashed lines in a and b mark R=1. Points and bands are condition means and pointwise 95\% complete-trajectory bootstrap intervals (10,000 resamples). All other settings are as in figure~\ref{fig:1}.}
\label{fig:2}
\end{figure}

\begin{figure}[!t]
\centering
\includegraphics[width=\linewidth,height=0.72\textheight,keepaspectratio]{figures/Figure_3.pdf}
\caption{\textbf{Public communication and contribution dynamics.} Each model block contains four time-course panels at R=0.60, 0.98, 1.00 and 1.02 (a--d) and a whole-run communication-associated difference map (e). Grey circles, blue circles and orange squares denote no communication, first-round public communication and every-round public communication. In e, CE denotes communication effect and pp denotes percentage points. CE is the whole-run mean under communication minus the corresponding no-communication mean, multiplied by 100. First and Every identify the two communication schedules. The five rows correspond to R=0.60, 0.90, 0.98, 1.00 and 1.02. Cell labels are rounded to whole percentage points. Exact differences and intervals are retained in source data. Means and contrasts use separate complete ensembles. Shaded time-course bands are pointwise 95\% intervals from 10,000 within-condition complete-trajectory bootstrap resamples. All games use N=6, T=10 and E=20, with fixed membership, full history and aggregate feedback.}
\label{fig:3}
\end{figure}

\begin{figure}[!t]
\centering
\includegraphics[width=\linewidth,height=0.72\textheight,keepaspectratio]{figures/Figure_4.pdf}
\caption{\textbf{Contribution-state composition across rounds and communication schedules.} Each model block contains no communication, first-round public communication and every-round public communication. At N=6, within-condition fractions are averaged equally over R=0.60, 0.90 and 0.98. Zero, partial and full contributions correspond to \(c_i(t)=0\), \(0<c_i(t)<E\) and \(c_i(t)=E\). Fractions sum to 100\% at every observed round. Areas connect the ten round values. No uncertainty bands are shown. Each condition includes 10--13 independent groups for DeepSeek and 30 for each other model.}
\label{fig:4}
\end{figure}

\begin{figure}[!t]
\centering
\includegraphics[width=\linewidth,height=0.72\textheight,keepaspectratio]{figures/Figure_5.pdf}
\caption{\textbf{Random termination and fixed-duration comparison.} Upper panels show individual completed group trajectories under no communication at R=0.90 and 1.02, $\delta$=0.90, N=6 and E=20. Each row is one group, ordered by realized duration. Coloured cells show observed round-specific contributions. A black tick marks the actual endpoint. Grey hatching marks unobserved periods after termination. Lower panels show each group's whole-run mean under fixed-ten-round or random termination. Horizontal jitter affects only categorical placement. Black diamonds and bars show equal-group means and 95\% intervals from 10,000 complete-trajectory bootstrap samples. The stopping rule is described in supplementary section~\ref{supp-sec:S1.5}.}
\label{fig:5}
\end{figure}

\begin{figure}[H]
\centering
\includegraphics[width=\linewidth,height=0.78\textheight,keepaspectratio]{figures/Figure_6.pdf}
\caption{\textbf{Group size and public communication.} Panels follow the model order in figure~\ref{fig:1}. Grey circles, blue circles and orange squares show no communication, first-round public communication and every-round public communication. Each point is the whole-run mean at N=4, 6, 8 or 12, averaging equally over R=0.60, 0.90 and 0.98. Error bars are pointwise 95\% percentile intervals from 10,000 within-condition bootstrap resamples of complete group trajectories. The comparison includes 216 conditions and 3,977 complete groups, with 9--30 groups per condition. Communication at N=4, 8 and 12 uses 10 groups per condition, except three Kimi conditions with nine completed groups. All games use T=10, E=20, announced termination, fixed membership, full history and aggregate feedback. R remains fixed as N changes. Dashed lines connect the size-specific estimates.}
\label{fig:6}
\end{figure}

\begin{figure}[H]
\centering
\includegraphics[width=\linewidth,height=0.78\textheight,keepaspectratio]{figures/Figure_7.pdf}
\caption{\textbf{Reasoning settings across incentives and group sizes.} DeepSeek occupies panels a1--d1, Claude a2--d2 and GPT a3--d3. Columns show N=4, 6, 8 and 12. Grey circles and blue squares identify reasoning disabled and enabled, respectively. Points give whole-run condition means across all 11 R values at their actual numerical spacing. Shading shows the original pointwise 95\% percentile intervals from 10,000 complete-trajectory bootstrap resamples. Vertical dashed lines mark R=1. The figure contains 264 conditions and 6,215 complete groups, with 10--30 groups per condition. All games use no communication, T=10, E=20, fixed membership, full history and aggregate feedback. The configuration pairs are specified in supplementary section~\ref{supp-sec:S1.1}.}
\label{fig:7}
\end{figure}

\begin{figure}[H]
\centering
\includegraphics[width=\linewidth]{figures/Figure_S1.pdf}
\caption{\textbf{Communication contrasts across outcome windows.} Each panel follows the three contrasts for one model across whole-run, late-stage and final-round contribution. Estimates average equally over R=0.60, 0.90 and 0.98 at N=6, T=10 and E=20. Blue circles denote first-round communication minus no communication, orange squares denote every-round communication minus no communication, and grey triangles denote every-round minus first-round communication. Whiskers retain the pointwise 95\% percentile intervals from 10,000 within-condition complete-group bootstrap resamples; communication arms are sampled independently, with shared control draws reused across related contrasts. Whole denotes rounds 1--10, late denotes rounds 8--10, and final denotes round 10. Dashed connectors link these overlapping summary windows; they are not round-by-round trajectories. Small horizontal offsets separate the three estimates at each window. Differences are in percentage points.}
\label{fig:S1}
\end{figure}

\begin{figure}[H]
\centering
\includegraphics[width=\linewidth]{figures/Figure_S2.pdf}
\caption{\textbf{Distributions of complete-group outcomes under communication.} Each model block pairs a frequency strip (left) with an empirical cumulative distribution (right) of whole-run contribution at R=0.90, N=6, T=10 and E=20. Grey, blue and orange denote no communication (None), first-round communication (First) and every-round communication (Every). Frequency markers use five-percentage-point bins, with area proportional to the number of groups; labels identify counts of three or more. Bins are left-closed and right-open, except that the final bin includes 100\%. Vertical displacement separates neighbouring markers. Black marks below each strip show the median and interquartile range calculated from exact outcomes. The step curves use exact unbinned outcomes and give the percentage of groups with contribution at or below the horizontal-axis value. All 483 groups are included: DeepSeek contributes 13/10/10 groups in None/First/Every order, and the other models contribute 30 per schedule. No distribution smoothing is applied.}
\label{fig:S2}
\end{figure}

\begin{figure}[H]
\centering
\includegraphics[width=\linewidth]{figures/Figure_S3.pdf}
\caption{\textbf{DeepSeek responses to task specification.} The five columns separate R=0.60, 0.90, 0.98, 1.00 and 1.02. Panels a--e show absolute whole-run contribution under changes in announced horizon (shaded rows) and feedback/history (remaining rows). Horizon values are T=1, 5, 10, 20 and 50. No feedback/history removes both; Aggregate and Individual specify the disclosed contributions, and 1, 3 or all specify the retained history. Panels f--j show matched whole-run differences for endowment, action resolution and membership. Baselines are E=20, step/E=0.05 and fixed membership, respectively. Points and whiskers retain all original means and pointwise 95\% percentile intervals from 10,000 complete-trajectory bootstrap resamples. Absolute contributions and contrasts have separate shared horizontal scales. All conditions use N=6 and no communication; E=20 and T=10 apply except where varied. Random regrouping uses a 24-participant pool, whose complete trajectory is the independent resampling unit.}
\label{fig:S3}
\end{figure}
